\documentclass{report}
\usepackage{amsmath,amssymb}
\setlength{\parindent}{0mm}
\setlength{\parskip}{1em}
\begin{document}
\begin{center}
\rule{6in}{1pt} \
{\large Max la Cour Christensen \\
{\bf Nonlinear Multigrid for Reservoir Simulation}}

Viggo Jarlsvej 2 \\ 207 \\ 2830 Virum \\ Denmark
\\
{\tt max@maxlacour.com}\\
Klaus Langgren Eskildsen\\
Allan P. Engsig-Karup\\
Mark Wakefield\end{center}

In the pursuit of higher resolution simulation models that use all
seismic, geological, and dynamic reservoir data - and to make use of
modern parallel computing architectures - we consider alternative
numerical methods to solve the system of equations governing subsurface
porous media flow.

The subject of this work is a thorough investigation of the application
of nonlinear multigrid techniques, specifically the Full Approximation
Scheme (FAS), for simulation of subsurface multiphase porous media flow.
The main motivation for addressing this topic is a need for higher
resolution and efficient simulations leading to better decision making in
the production of oil and gas. Higher resolution simulations require
efficient utilization of many-core parallel architectures. Current
numerical methods employed in industrial reservoir simulators are memory
intensive and not readily scalable on large-scale distributed systems and
modern many-core architectures such as GPUs or Intel MICs.

In a first step, we investigate alternative numerical methods to
establish algorithmic performance in serial computations. The nonlinear
multigrid technique FAS uses local linearization, which allows for local
components suitable for parallel implementation. Furthermore, FAS is a
memory lean algorithm.
By using FAS, we avoid having to assemble the Jacobian on the finest
grid, which results in major memory savings. To our knowledge, very
little work is published on FAS for reservoir simulation. Previous work
by Molenaar considers the application of FAS on a simple 2D immiscible
two-phase no gravity homogeneous example. To our knowledge, FAS has not
been applied successfully to more complicated heterogeneous reservoir
problems.

Two reservoir simulators have been implemented in C++ in serial. The
first simulator is based on conventional techniques with global
linearization in Newton�s method and state-of-the-art choice of methods
for the linear solver. Specifically, we use the Krylov subspace method
FGMRES preconditioned with the Constrained Pressure Residual (CPR)
method. CPR is a two-stage preconditioner based on a decomposition of the
full system into a pressure system, for which the solution is then used
to correct the full system. The intermediate pressure system is solved by
an algebraic multigrid method. The second stage of CPR is done by
applying ILU(0) to the corrected full system.

The second simulator is based on the nonlinear multigrid method FAS. Both
simulators solve the same system of PDEs governing 3D three-phase flow of
oil, water and gas in a subsurface porous medium taking into account
gravitational effects. The same discretization techniques are used for
both simulators. For spatial discretization, the Finite Volume Method is
used and for temporal integration, the backward Euler method is used.
This enables fair comparisons between the conventional methods and FAS.

The two reservoir simulators have been tested extensively to compare the
nonlinear multigrid approach FAS with the conventional techniques applied
in modern reservoir simulation. The test models range from homogeneous to
highly heterogeneous models. It has been demonstrated that, without loss
of robustness, FAS outperforms the conventional techniques in terms of
algorithmic and numerical efficiency for the model equations considered.
Furthermore, memory comparisons have been carried out, which show that
FAS provides a significant memory reduction in comparison with
conventional techniques. This memory reduction is an attractive feature,
which enables higher resolution simulation for the before mentioned
modern many-core architectures.


\end{document}
